% =====================================================================
%  Appendix C: Matrix Algebra Quick Reference
% =====================================================================

\chapter{Matrix Algebra Quick Reference}
\label{app:matrix-algebra-reference}

Matrices are the bookkeeping language for linear maps.  In this book they appear as
linear transformations, Jacobian matrices, Hessian matrices, coordinate changes, and
quadratic forms.

This appendix is a reference.  It is meant to be scanned.

\[
\begin{array}{c|c}
\text{Matrix idea} & \text{Where it appears in the book} \\ \hline
\text{matrix multiplication} & \text{composition of linear maps}\\[1mm]
\text{row reduction} & \text{solving linear systems}\\[1mm]
\text{determinant} & \text{area, volume, orientation, Jacobians}\\[1mm]
\text{inverse matrix} & \text{undoing a linear map}\\[1mm]
\text{eigenvalues} & \text{principal directions, quadratic forms}\\[1mm]
\text{quadratic forms} & \text{Hessian test, conics, local shape}
\end{array}
\]

% ---------------------------------------------------------------------
\section{Matrix operations}
\label{app:C.1}

\subsection*{Matrix notation}

An \(m\times n\) matrix has \(m\) rows and \(n\) columns:
\[
\boxed{
    A=
    \begin{pmatrix}
    a_{11} & a_{12} & \cdots & a_{1n}\\
    a_{21} & a_{22} & \cdots & a_{2n}\\
    \vdots & \vdots & \ddots & \vdots\\
    a_{m1} & a_{m2} & \cdots & a_{mn}
    \end{pmatrix}.
}
\]

The entry in row \(i\), column \(j\) is
\[
\boxed{
    a_{ij}.
}
\]

A column vector in \(\mathbb R^n\) is an \(n\times1\) matrix:
\[
\boxed{
    \mathbf x=
    \begin{pmatrix}
    x_1\\
    x_2\\
    \vdots\\
    x_n
    \end{pmatrix}.
}
\]

\subsection*{Addition}

Matrices can be added only when they have the same size.

\[
\boxed{
    (A+B)_{ij}=a_{ij}+b_{ij}.
}
\]

\subsection*{Scalar multiplication}

\[
\boxed{
    (cA)_{ij}=c\,a_{ij}.
}
\]

\subsection*{Matrix-vector multiplication}

If \(A\) is \(m\times n\) and \(\mathbf x\in\mathbb R^n\), then
\[
\boxed{
    A\mathbf x\in\mathbb R^m.
}
\]

The \(i\)th component is
\[
\boxed{
    (A\mathbf x)_i
    =
    a_{i1}x_1+a_{i2}x_2+\cdots+a_{in}x_n.
}
\]

Equivalently,
\[
\boxed{
    A\mathbf x
    =
    x_1\mathbf a_1+x_2\mathbf a_2+\cdots+x_n\mathbf a_n,
}
\]
where \(\mathbf a_1,\ldots,\mathbf a_n\) are the columns of \(A\).

\subsection*{Matrix multiplication}

If \(A\) is \(m\times n\) and \(B\) is \(n\times p\), then
\[
\boxed{
    AB\text{ is }m\times p.
}
\]

The entry in row \(i\), column \(j\) is
\[
\boxed{
    (AB)_{ij}
    =
    \sum_{k=1}^{n}a_{ik}b_{kj}.
}
\]

Matrix multiplication means composition of linear maps:
\[
\boxed{
    A(B\mathbf x)=(AB)\mathbf x.
}
\]

\subsection*{Identity matrix}

The \(n\times n\) identity matrix is
\[
\boxed{
    I_n=
    \begin{pmatrix}
    1&0&\cdots&0\\
    0&1&\cdots&0\\
    \vdots&\vdots&\ddots&\vdots\\
    0&0&\cdots&1
    \end{pmatrix}.
}
\]

It satisfies
\[
\boxed{
    I_n\mathbf x=\mathbf x.
}
\]

If the sizes match,
\[
\boxed{
    AI=IA=A.
}
\]

\subsection*{Zero matrix}

The zero matrix has all entries equal to \(0\).  It satisfies
\[
\boxed{
    A+0=A.
}
\]

If the sizes match,
\[
\boxed{
    A0=0,
    \qquad
    0A=0.
}
\]

\subsection*{Transpose}

The transpose of \(A\) is \(A^T\).  Rows become columns.

\[
\boxed{
    (A^T)_{ij}=a_{ji}.
}
\]

Rules:
\[
\boxed{
    (A+B)^T=A^T+B^T
}
\]

\[
\boxed{
    (cA)^T=cA^T
}
\]

\[
\boxed{
    (AB)^T=B^T A^T
}
\]

\[
\boxed{
    (A^T)^T=A.
}
\]

\subsection*{Symmetric matrix}

A square matrix is symmetric if
\[
\boxed{
    A^T=A.
}
\]

Symmetric matrices are especially important for Hessians and quadratic forms.

\subsection*{Common warnings}

Matrix multiplication is usually not commutative:
\[
\boxed{
    AB\ne BA
    \quad\text{in general}.
}
\]

The product \(AB\) may exist even when \(BA\) does not.

If both \(AB\) and \(BA\) exist, they may have different sizes.

Always check dimensions before multiplying.

% ---------------------------------------------------------------------
\section{Row reduction}
\label{app:C.2}

\subsection*{Linear systems}

A system of linear equations can be written as
\[
\boxed{
    A\mathbf x=\mathbf b.
}
\]

The augmented matrix is
\[
\boxed{
    [A\mid \mathbf b].
}
\]

\subsection*{Elementary row operations}

The allowed row operations are:

\[
\begin{array}{ll}
1. & \text{Swap two rows.}\\[1mm]
2. & \text{Multiply a row by a nonzero scalar.}\\[1mm]
3. & \text{Add a multiple of one row to another row.}
\end{array}
\]

These operations do not change the solution set of the system.

\subsection*{Row echelon form}

A matrix is in row echelon form if:

\[
\begin{array}{ll}
1. & \text{All zero rows are at the bottom.}\\[1mm]
2. & \text{Each leading entry is to the right of the leading entry above it.}\\[1mm]
3. & \text{All entries below a leading entry are zero.}
\end{array}
\]

\subsection*{Reduced row echelon form}

A matrix is in reduced row echelon form if it is in row echelon form and:

\[
\begin{array}{ll}
1. & \text{Each leading entry is }1.\\[1mm]
2. & \text{Each leading }1\text{ is the only nonzero entry in its column.}
\end{array}
\]

\subsection*{Pivot}

A pivot is a leading nonzero entry in a row.

A pivot column is a column containing a pivot.

Variables corresponding to pivot columns are pivot variables.

Variables corresponding to non-pivot columns are free variables.

\subsection*{Rank}

The rank of a matrix is the number of pivots:
\[
\boxed{
    \operatorname{rank}(A)=\text{number of pivot columns in }A.
}
\]

\subsection*{Consistency test}

For the system
\[
    A\mathbf x=\mathbf b,
\]
row reduce
\[
    [A\mid \mathbf b].
\]

If a row has the form
\[
\boxed{
    0\quad 0\quad \cdots\quad 0\mid c
}
\]
with
\[
    c\ne0,
\]
then the system has no solution.

Otherwise, the system is consistent.

\subsection*{Number of solutions}

For a consistent system:

\[
\begin{array}{c|c}
\text{Situation} & \text{Solution type} \\ \hline
\text{pivot in every variable column} & \text{unique solution}\\[1mm]
\text{at least one free variable} & \text{infinitely many solutions}
\end{array}
\]

\subsection*{Homogeneous systems}

A homogeneous system has the form
\[
\boxed{
    A\mathbf x=\mathbf 0.
}
\]

It is always consistent because
\[
    \mathbf x=\mathbf 0
\]
is always a solution.

This is called the trivial solution.

A homogeneous system has a nontrivial solution exactly when there is at least one free
variable.

\subsection*{Null space}

The null space of \(A\) is
\[
\boxed{
    \operatorname{Null}(A)
    =
    \{\mathbf x:A\mathbf x=\mathbf 0\}.
}
\]

It is the set of vectors that \(A\) sends to zero.

\subsection*{Column space}

The column space of \(A\) is
\[
\boxed{
    \operatorname{Col}(A)
    =
    \text{span of the columns of }A.
}
\]

The equation
\[
    A\mathbf x=\mathbf b
\]
has a solution exactly when
\[
\boxed{
    \mathbf b\in\operatorname{Col}(A).
}
\]

\subsection*{Common warnings}

Row operations preserve the solution set of a linear system, but they do not usually
preserve determinants.

Do not row reduce \(A\) alone if the system is
\[
    A\mathbf x=\mathbf b.
\]
Row reduce the augmented matrix
\[
    [A\mid \mathbf b].
\]

Free variables are not errors.  They describe infinitely many solutions.

% ---------------------------------------------------------------------
\section{Determinants}
\label{app:C.3}

\subsection*{Only square matrices have determinants}

The determinant is defined for square matrices:
\[
\boxed{
    A\text{ is }n\times n
    \quad\Longrightarrow\quad
    \det A\text{ is a number}.
}
\]

\subsection*{\(2\times2\) determinant}

\[
\boxed{
    \det
    \begin{pmatrix}
    a&b\\
    c&d
    \end{pmatrix}
    =
    ad-bc.
}
\]

Geometric meaning:

\[
\boxed{
    |\det A|
    =
    \text{area scale factor of the linear map }A.
}
\]

The sign records orientation.

\subsection*{\(3\times3\) determinant}

\[
\boxed{
\det
\begin{pmatrix}
a&b&c\\
d&e&f\\
g&h&i
\end{pmatrix}
=
a(ei-fh)-b(di-fg)+c(dh-eg).
}
\]

Geometric meaning:

\[
\boxed{
    |\det A|
    =
    \text{volume scale factor of the linear map }A.
}
\]

The sign records orientation.

\subsection*{Determinant and orientation}

\[
\boxed{
    \det A>0
}
\]
means orientation is preserved.

\[
\boxed{
    \det A<0
}
\]
means orientation is reversed.

\[
\boxed{
    \det A=0
}
\]
means area or volume collapses.

\subsection*{Singular matrices}

For a square matrix \(A\),
\[
\boxed{
    \det A=0
}
\]
means \(A\) is singular.

Equivalently, the linear map collapses some nonzero direction to zero.

\subsection*{Determinant properties}

\[
\boxed{
    \det(I)=1
}
\]

\[
\boxed{
    \det(A^T)=\det(A)
}
\]

\[
\boxed{
    \det(AB)=\det(A)\det(B)
}
\]

If \(A\) is invertible,
\[
\boxed{
    \det(A^{-1})=\frac1{\det(A)}.
}
\]

If \(A\) is triangular, then
\[
\boxed{
    \det(A)=\text{product of diagonal entries}.
}
\]

\subsection*{Effect of row operations}

Swapping two rows changes the sign of the determinant.

Multiplying one row by \(c\) multiplies the determinant by \(c\).

Adding a multiple of one row to another row does not change the determinant.

\subsection*{Determinant and linear independence}

For an \(n\times n\) matrix \(A\),
\[
\boxed{
    \det A\ne0
}
\]
if and only if the columns of \(A\) are linearly independent.

\[
\boxed{
    \det A=0
}
\]
if and only if the columns of \(A\) are linearly dependent.

\subsection*{Determinant and change of variables}

For a coordinate map
\[
    (u,v)\mapsto (x(u,v),y(u,v)),
\]
the Jacobian determinant is
\[
\boxed{
    \frac{\partial(x,y)}{\partial(u,v)}
    =
    \det
    \begin{pmatrix}
    x_u & x_v\\
    y_u & y_v
    \end{pmatrix}.
}
\]

For ordinary area,
\[
\boxed{
    dA
    =
    \left|
    \frac{\partial(x,y)}{\partial(u,v)}
    \right|
    \,du\,dv.
}
\]

For oriented area,
\[
\boxed{
    dx\wedge dy
    =
    \frac{\partial(x,y)}{\partial(u,v)}
    \,du\wedge dv.
}
\]

\subsection*{Common warnings}

The determinant is not linear in the whole matrix:
\[
    \det(A+B)\ne\det(A)+\det(B)
\]
in general.

A determinant can be negative.  Ordinary area and volume use \(|\det A|\).

If row reduction uses row swaps or row scaling, track how the determinant changes.

% ---------------------------------------------------------------------
\section{Inverse matrices}
\label{app:C.4}

\subsection*{Definition}

A square matrix \(A\) is invertible if there is a matrix \(A^{-1}\) such that
\[
\boxed{
    A^{-1}A=AA^{-1}=I.
}
\]

The matrix \(A^{-1}\) is the inverse of \(A\).

\subsection*{Meaning}

If
\[
    \mathbf y=A\mathbf x,
\]
then
\[
\boxed{
    \mathbf x=A^{-1}\mathbf y.
}
\]

The inverse matrix undoes the linear map \(A\).

\subsection*{\(2\times2\) inverse formula}

If
\[
    A=
    \begin{pmatrix}
    a&b\\
    c&d
    \end{pmatrix},
\]
then
\[
\boxed{
    A^{-1}
    =
    \frac{1}{ad-bc}
    \begin{pmatrix}
    d&-b\\
    -c&a
    \end{pmatrix}
}
\]
provided
\[
\boxed{
    ad-bc\ne0.
}
\]

\subsection*{Invertibility test}

For an \(n\times n\) matrix \(A\), the following statements are equivalent:

\[
\begin{array}{ll}
1. & A\text{ is invertible.}\\[1mm]
2. & \det A\ne0.\\[1mm]
3. & A\mathbf x=\mathbf 0\text{ has only the trivial solution.}\\[1mm]
4. & A\mathbf x=\mathbf b\text{ has a unique solution for every }\mathbf b.\\[1mm]
5. & \text{The columns of }A\text{ are linearly independent.}\\[1mm]
6. & \text{The columns of }A\text{ span }\mathbb R^n.\\[1mm]
7. & \operatorname{rref}(A)=I.
\end{array}
\]

\subsection*{Finding an inverse by row reduction}

To find \(A^{-1}\), row reduce
\[
\boxed{
    [A\mid I].
}
\]

If
\[
    [A\mid I]
    \longrightarrow
    [I\mid B],
\]
then
\[
\boxed{
    B=A^{-1}.
}
\]

If the left side cannot be reduced to \(I\), then \(A\) is not invertible.

\subsection*{Solving systems with inverses}

If \(A\) is invertible, then
\[
\boxed{
    A\mathbf x=\mathbf b
    \quad\Longrightarrow\quad
    \mathbf x=A^{-1}\mathbf b.
}
\]

\subsection*{Inverse rules}

If \(A\) and \(B\) are invertible, then
\[
\boxed{
    (AB)^{-1}=B^{-1}A^{-1}.
}
\]

\[
\boxed{
    (A^{-1})^{-1}=A.
}
\]

\[
\boxed{
    (A^T)^{-1}=(A^{-1})^T.
}
\]

\[
\boxed{
    (cA)^{-1}=\frac1c A^{-1}
    \qquad (c\ne0).
}
\]

\subsection*{Common warnings}

The inverse of a product reverses the order:
\[
    (AB)^{-1}=B^{-1}A^{-1}.
\]

Not every square matrix has an inverse.

Non-square matrices do not have two-sided inverses.

In practice, row reduction is usually better than explicitly computing \(A^{-1}\) to solve
a system.

% ---------------------------------------------------------------------
\section{Eigenvalues and eigenvectors}
\label{app:C.5}

\subsection*{Definition}

Let \(A\) be an \(n\times n\) matrix.

A nonzero vector \(\mathbf v\) is an eigenvector of \(A\) if
\[
\boxed{
    A\mathbf v=\lambda \mathbf v
}
\]
for some scalar \(\lambda\).

The scalar \(\lambda\) is the eigenvalue.

\subsection*{Meaning}

An eigenvector is a direction that \(A\) does not turn.

The matrix \(A\) only stretches, shrinks, or reverses that direction.

\[
\begin{array}{c|c}
\lambda & \text{Effect on eigenvector direction} \\ \hline
\lambda>1 & \text{stretch}\\
0<\lambda<1 & \text{shrink}\\
\lambda<0 & \text{reverse direction and scale}\\
\lambda=0 & \text{collapse to zero}
\end{array}
\]

\subsection*{Finding eigenvalues}

Start with
\[
    A\mathbf v=\lambda\mathbf v.
\]
Rewrite:
\[
    (A-\lambda I)\mathbf v=\mathbf 0.
\]
For a nonzero solution \(\mathbf v\), the matrix \(A-\lambda I\) must be singular:
\[
\boxed{
    \det(A-\lambda I)=0.
}
\]

This equation is the characteristic equation.

\subsection*{Characteristic polynomial}

\[
\boxed{
    p(\lambda)=\det(A-\lambda I).
}
\]

Eigenvalues are the roots of
\[
\boxed{
    p(\lambda)=0.
}
\]

\subsection*{\(2\times2\) shortcut}

If
\[
    A=
    \begin{pmatrix}
    a&b\\
    c&d
    \end{pmatrix},
\]
then
\[
\boxed{
    \operatorname{tr}(A)=a+d
}
\]
and
\[
\boxed{
    \det(A)=ad-bc.
}
\]

The characteristic equation can be written as
\[
\boxed{
    \lambda^2-\operatorname{tr}(A)\lambda+\det(A)=0.
}
\]

\subsection*{Finding eigenvectors}

For each eigenvalue \(\lambda\), solve
\[
\boxed{
    (A-\lambda I)\mathbf v=\mathbf 0.
}
\]

The solution space is the eigenspace for \(\lambda\).

\[
\boxed{
    E_\lambda=\operatorname{Null}(A-\lambda I).
}
\]

Remember:
\[
\boxed{
    \mathbf v\ne\mathbf 0.
}
\]
The zero vector is never called an eigenvector.

\subsection*{Diagonal matrices}

If
\[
    D=
    \begin{pmatrix}
    \lambda_1&0&\cdots&0\\
    0&\lambda_2&\cdots&0\\
    \vdots&\vdots&\ddots&\vdots\\
    0&0&\cdots&\lambda_n
    \end{pmatrix},
\]
then the standard basis vectors are eigenvectors:
\[
\boxed{
    D\mathbf e_i=\lambda_i\mathbf e_i.
}
\]

\subsection*{Diagonalization}

If \(A\) has \(n\) linearly independent eigenvectors
\[
    \mathbf v_1,\ldots,\mathbf v_n,
\]
then put them into a matrix
\[
\boxed{
    P=
    \begin{pmatrix}
    |&|&&|\\
    \mathbf v_1&\mathbf v_2&\cdots&\mathbf v_n\\
    |&|&&|
    \end{pmatrix}.
}
\]

Let
\[
\boxed{
    D=
    \begin{pmatrix}
    \lambda_1&0&\cdots&0\\
    0&\lambda_2&\cdots&0\\
    \vdots&\vdots&\ddots&\vdots\\
    0&0&\cdots&\lambda_n
    \end{pmatrix}.
}
\]

Then
\[
\boxed{
    A=PDP^{-1}.
}
\]

\subsection*{Powers}

If
\[
    A=PDP^{-1},
\]
then
\[
\boxed{
    A^k=PD^kP^{-1}.
}
\]

For a diagonal matrix,
\[
\boxed{
    D^k=
    \begin{pmatrix}
    \lambda_1^k&0&\cdots&0\\
    0&\lambda_2^k&\cdots&0\\
    \vdots&\vdots&\ddots&\vdots\\
    0&0&\cdots&\lambda_n^k
    \end{pmatrix}.
}
\]

\subsection*{Symmetric matrices}

If
\[
    A^T=A,
\]
then:

\[
\begin{array}{ll}
1. & \text{All eigenvalues of }A\text{ are real.}\\[1mm]
2. & \text{Eigenvectors for different eigenvalues are orthogonal.}\\[1mm]
3. & A\text{ has an orthonormal eigenbasis.}
\end{array}
\]

Thus a symmetric matrix can be written as
\[
\boxed{
    A=QDQ^T,
}
\]
where \(Q\) is orthogonal:
\[
\boxed{
    Q^TQ=I.
}
\]

Symmetric matrices are the clean case.  Hessian matrices are symmetric when mixed
partials agree.

\subsection*{Common warnings}

Eigenvectors must be nonzero.

Eigenvalues can be repeated.

A matrix may fail to have enough eigenvectors to diagonalize.

Eigenvalues are defined only for square matrices.

For real matrices, complex eigenvalues may occur unless the matrix is symmetric.

% ---------------------------------------------------------------------
\section{Quadratic forms}
\label{app:C.6}

\subsection*{Definition}

A quadratic form is a function built from degree-two terms.

In matrix form:
\[
\boxed{
    q(\mathbf x)=\mathbf x^T A\mathbf x.
}
\]

Usually \(A\) is taken to be symmetric.

\subsection*{Two-variable form}

For
\[
    \mathbf x=
    \begin{pmatrix}
    x\\
    y
    \end{pmatrix},
\]
and
\[
    A=
    \begin{pmatrix}
    a&b\\
    b&c
    \end{pmatrix},
\]
we get
\[
\boxed{
    \mathbf x^T A\mathbf x
    =
    ax^2+2bxy+cy^2.
}
\]

So the form
\[
    ax^2+2bxy+cy^2
\]
corresponds to the matrix
\[
\boxed{
    \begin{pmatrix}
    a&b\\
    b&c
    \end{pmatrix}.
}
\]

\subsection*{Symmetric part}

For any square matrix \(A\),
\[
\boxed{
    \mathbf x^T A\mathbf x
    =
    \mathbf x^T\left(\frac{A+A^T}{2}\right)\mathbf x.
}
\]

So only the symmetric part of \(A\) matters for a quadratic form.

\subsection*{Positive definite}

A quadratic form is positive definite if
\[
\boxed{
    q(\mathbf x)>0
    \quad\text{for every }\mathbf x\ne\mathbf 0.
}
\]

A symmetric matrix \(A\) is positive definite if
\[
\boxed{
    \mathbf x^T A\mathbf x>0
    \quad\text{for every }\mathbf x\ne\mathbf 0.
}
\]

\subsection*{Negative definite}

A quadratic form is negative definite if
\[
\boxed{
    q(\mathbf x)<0
    \quad\text{for every }\mathbf x\ne\mathbf 0.
}
\]

\subsection*{Indefinite}

A quadratic form is indefinite if it takes both positive and negative values.

\[
\boxed{
    q(\mathbf u)>0
    \quad\text{and}\quad
    q(\mathbf v)<0
}
\]
for some vectors \(\mathbf u,\mathbf v\).

\subsection*{Semidefinite}

Positive semidefinite:
\[
\boxed{
    q(\mathbf x)\ge0
    \quad\text{for every }\mathbf x.
}
\]

Negative semidefinite:
\[
\boxed{
    q(\mathbf x)\le0
    \quad\text{for every }\mathbf x.
}
\]

Semidefinite forms can be zero on nonzero vectors.

\subsection*{\(2\times2\) test}

Let
\[
    A=
    \begin{pmatrix}
    a&b\\
    b&c
    \end{pmatrix}
\]
and
\[
    D=\det A=ac-b^2.
\]

\[
\begin{array}{c|c}
\text{Condition} & \text{Type} \\ \hline
D>0,\ a>0 & \text{positive definite}\\[1mm]
D>0,\ a<0 & \text{negative definite}\\[1mm]
D<0 & \text{indefinite}\\[1mm]
D=0 & \text{inconclusive from this test}
\end{array}
\]

\subsection*{Eigenvalue test}

For a symmetric matrix \(A\):

\[
\begin{array}{c|c}
\text{Eigenvalues of }A & \text{Type of quadratic form} \\ \hline
\text{all positive} & \text{positive definite}\\[1mm]
\text{all negative} & \text{negative definite}\\[1mm]
\text{positive and negative} & \text{indefinite}\\[1mm]
\text{all nonnegative, at least one zero} & \text{positive semidefinite}\\[1mm]
\text{all nonpositive, at least one zero} & \text{negative semidefinite}
\end{array}
\]

\subsection*{Diagonalization of a quadratic form}

If \(A\) is symmetric, then
\[
\boxed{
    A=QDQ^T,
}
\]
where \(Q\) is orthogonal and \(D\) is diagonal.

Let
\[
    \mathbf y=Q^T\mathbf x.
\]
Then
\[
\boxed{
    \mathbf x^T A\mathbf x
    =
    \mathbf y^T D\mathbf y
    =
    \lambda_1 y_1^2+\cdots+\lambda_n y_n^2.
}
\]

The eigenvalues show the principal curvatures of the quadratic form.

\subsection*{Connection to Hessians}

For a function \(f(x,y)\), the Hessian matrix is
\[
\boxed{
    H_f=
    \begin{pmatrix}
    f_{xx} & f_{xy}\\
    f_{yx} & f_{yy}
    \end{pmatrix}.
}
\]

When mixed partials agree,
\[
    f_{xy}=f_{yx},
\]
so \(H_f\) is symmetric.

Near a critical point \(\mathbf p\), the quadratic approximation is
\[
\boxed{
    f(\mathbf p+\mathbf h)
    \approx
    f(\mathbf p)
    +
    \frac12\mathbf h^T H_f(\mathbf p)\mathbf h.
}
\]

\subsection*{Second derivative test in two variables}

Suppose
\[
    \nabla f(a,b)=\mathbf 0.
\]

Let
\[
\boxed{
    D=f_{xx}(a,b)f_{yy}(a,b)-f_{xy}(a,b)^2.
}
\]

Then:

\[
\begin{array}{c|c}
\text{Condition} & \text{Conclusion} \\ \hline
D>0,\ f_{xx}(a,b)>0 & \text{local minimum}\\[1mm]
D>0,\ f_{xx}(a,b)<0 & \text{local maximum}\\[1mm]
D<0 & \text{saddle point}\\[1mm]
D=0 & \text{test inconclusive}
\end{array}
\]

\subsection*{Common warnings}

A zero determinant in the Hessian test does not mean there is no maximum or minimum.
It means the quadratic test does not decide.

The matrix for
\[
    ax^2+bxy+cy^2
\]
is
\[
    \begin{pmatrix}
    a & b/2\\
    b/2 & c
    \end{pmatrix},
\]
not
\[
    \begin{pmatrix}
    a & b\\
    b & c
    \end{pmatrix}.
\]

The form
\[
    \mathbf x^T A\mathbf x
\]
uses a symmetric matrix when reading geometry.

% ---------------------------------------------------------------------
\section*{One-page formula summary}

\[
\boxed{
    (A\mathbf x)_i
    =
    \sum_{j=1}^n a_{ij}x_j
}
\]

\[
\boxed{
    (AB)_{ij}
    =
    \sum_k a_{ik}b_{kj}
}
\]

\[
\boxed{
    (AB)^T=B^T A^T
}
\]

\[
\boxed{
    \det
    \begin{pmatrix}
    a&b\\
    c&d
    \end{pmatrix}
    =
    ad-bc
}
\]

\[
\boxed{
    \det(AB)=\det(A)\det(B)
}
\]

\[
\boxed{
    A^{-1}
    =
    \frac{1}{ad-bc}
    \begin{pmatrix}
    d&-b\\
    -c&a
    \end{pmatrix}
}
\]
for
\[
    A=
    \begin{pmatrix}
    a&b\\
    c&d
    \end{pmatrix},
    \qquad
    ad-bc\ne0.
\]

\[
\boxed{
    A\text{ invertible}
    \Longleftrightarrow
    \det A\ne0
}
\]

\[
\boxed{
    A\mathbf v=\lambda\mathbf v
}
\]

\[
\boxed{
    \det(A-\lambda I)=0
}
\]

\[
\boxed{
    A=PDP^{-1}
}
\]
when \(A\) has enough independent eigenvectors.

\[
\boxed{
    q(\mathbf x)=\mathbf x^T A\mathbf x
}
\]

\[
\boxed{
    ax^2+2bxy+cy^2
    =
    \begin{pmatrix}x&y\end{pmatrix}
    \begin{pmatrix}
    a&b\\
    b&c
    \end{pmatrix}
    \begin{pmatrix}
    x\\
    y
    \end{pmatrix}
}
\]

\[
\boxed{
    A\text{ positive definite}
    \Longleftrightarrow
    \mathbf x^T A\mathbf x>0
    \text{ for every }\mathbf x\ne\mathbf 0
}
\]

\[
\boxed{
    D=f_{xx}f_{yy}-f_{xy}^2
}
\]

\[
\boxed{
\begin{array}{c|c}
D>0,\ f_{xx}>0 & \text{local minimum}\\
D>0,\ f_{xx}<0 & \text{local maximum}\\
D<0 & \text{saddle}\\
D=0 & \text{inconclusive}
\end{array}
}
\]